\subsection{Real-World Problem Description}

Traffic congestion is a common problem in urban
areas, particularly at intersections controlled
by traffic signals. Poorly coordinated traffic
signals can cause vehicles to stop repeatedly,
increase waiting times, and contribute to
congestion along major roads.

One issue occurs when traffic signals release
vehicles from one intersection while another
intersection farther down the road is unable
to accommodate the incoming traffic.

This problem motivated the investigation of
traffic signal scheduling using a greedy
algorithm.

For this project, we consider a simplified
traffic signal scheduling problem.

A traffic management system is assumed to
provide a collection of candidate green-light
time windows for a particular traffic movement.

Each window represents a time interval during
which a green-light opportunity could be
scheduled.

The scheduling system must select a subset
of these windows such that no two selected
windows overlap.

The objective is to maximize the total
number of compatible green-light windows.

This simplified model assumes that all
candidate windows have equal importance,
each window is indivisible, and the safety
and feasibility of every candidate window
have already been verified.

Although real-world traffic signal
coordination involves additional factors,
including traffic volume, vehicle speeds,
pedestrian crossings, and intersection
capacity, this abstraction allows us
to study an optimization problem that
can be solved optimally using a
greedy algorithm.

Traffic signal coordination is used to synchronize
intersections to improve traffic progression and
reduce unnecessary stops and delays \cite{fhwa2008}.

\subsection{Mathematical Abstraction}

The traffic signal scheduling problem
is modeled as a set of time intervals.

Let

\[
W = \{w_1, w_2, \ldots, w_n\}
\]

represent the set of candidate green-light
windows.

Each window is represented as

\[
w_i = (s_i, f_i)
\]

where:

\begin{itemize}
    \item $s_i$ represents the starting
    time of window $i$.

    \item $f_i$ represents the finishing
    time of window $i$.

    \item $n$ represents the total number
    of candidate windows.
\end{itemize}

Every window must satisfy

\[
s_i < f_i.
\]

Two windows are compatible when they
do not overlap.

For two windows $w_i$ and $w_j$,
where $w_i$ is scheduled before $w_j$,
compatibility requires

\[
f_i \leq s_j.
\]

This means that the next window may
begin only after the previous
selected window has finished.

The goal is to find a subset

\[
S \subseteq W
\]

that maximizes the number of
selected windows:

\[
\max |S|
\]

subject to the constraint that
all selected windows are
mutually compatible.

This is a maximum-cardinality
interval scheduling problem.

The greedy algorithm solves this
problem by sorting candidate
windows according to their
finishing times and repeatedly
selecting the earliest-finishing
window that does not conflict
with previously selected windows.

The result is an optimal subset
containing the maximum possible
number of compatible green-light
windows under the stated assumptions.